\documentclass[11pt,a4paper]{article}

% ─────────────────────────────────────────────────────────────
% Packages
% ─────────────────────────────────────────────────────────────
\usepackage[utf8]{inputenc}
\usepackage[T1]{fontenc}
\usepackage[french]{babel}
\usepackage{lmodern}
\usepackage{geometry}
\usepackage{amsmath,amssymb,amsfonts}
\usepackage{bm}
\usepackage{graphicx}
\usepackage{booktabs}
\usepackage{hyperref}
\usepackage{physics}
\usepackage{caption}
\usepackage{subcaption}
\usepackage{xcolor}
\usepackage[most]{tcolorbox}

\geometry{margin=2.5cm}

\hypersetup{
    colorlinks=true,
    linkcolor=blue,
    citecolor=blue,
    urlcolor=blue
}

% ─────────────────────────────────────────────────────────────
% Macros
% ─────────────────────────────────────────────────────────────
\newcommand{\bup}{\mathrm{BuP}}
\newcommand{\ent}{\mathrm{ent}}
\newcommand{\matter}{\mathrm{matter}}
\newcommand{\eff}{\mathrm{eff}}
\newcommand{\loc}{\mathrm{loc}}
\newcommand{\proxy}{\mathrm{proxy}}
\newcommand{\Omat}{\mathcal{O}_{\rm matter}}
\newcommand{\SbuP}{S_{\rm BuP}}
\newcommand{\W}{W_{ij}}

% ─────────────────────────────────────────────────────────────
% Title
% ─────────────────────────────────────────────────────────────
\title{
\textbf{Dynamique variationnelle des excitations BuP}\\[0.5em]
\large Métastabilité, phases internes et relaxation de la matière émergente}
\author{
Farid Hamdad\\
\small Independent Researcher\\
\small Bottom-Up Quantum Gravity Program
}
\date{2026}

\begin{document}

\maketitle

% ─────────────────────────────────────────────────────────────
\begin{abstract}
% ─────────────────────────────────────────────────────────────

Paper~6 a proposé une formulation pure du programme Bottom-Up Quantum Gravity
(BuP) sous la forme d'une équation fondamentale sur le réseau d'information
mutuelle,
\[
\mathcal{F}_{ij}[W]=0,
\]
ou, de manière variationnelle,
\[
\frac{\delta S_{\rm BuP}[W]}{\delta W_{ij}}=0.
\]
Dans cette formulation, la géométrie, la dimension spectrale et la matière ne
sont pas fondamentales : elles doivent émerger du réseau d'intrication
\[
W_{ij}=I(i:j).
\]

Le présent papier explore la dynamique induite par une action candidate
\(S_{\rm BuP}[W]\). Nous étudions comment une excitation localisée du réseau,
\[
W_{ij}=W_{ij}^{(0)}+\delta W_{ij}^{\rm loc},
\]
évolue sous relaxation variationnelle. Le caractère source d'une excitation
est mesuré par l'observable
\[
\mathcal{O}_{\rm matter}
=
\rho_{\rm Spearman}
\left(
T_{ij}^{\rm matter,proxy},
|\delta G_{ij}^{\rm proxy}|
\right),
\]
où
\[
T_{ij}^{\rm matter,proxy}=(\delta W_{ij})^2
\]
et
\[
G_{ij}^{\rm proxy}
=
-\kappa_{ij}+\frac12R_{\rm edge}.
\]

Le résultat principal est que \(\mathcal{O}_{\rm matter}\) n'est pas conservée
sous relaxation. Le seuil initial source/non-source est situé autour de
\[
\lambda_{\rm init}\simeq0.265.
\]
Après une relaxation intermédiaire de \(10\) à \(20\) pas, un plateau
métastable apparaît autour de
\[
\lambda_{\rm meta}\simeq0.876.
\]
Pour des relaxations plus longues, de \(40\) à \(320\) pas, le système rejoint
un régime tardif stable autour de
\[
\lambda_{\rm late}\simeq0.280.
\]
Un scan en taille finie (\(N=15,18,20,22,24\)) montre que cette dynamique
varie fortement avec \(N\) : elle est absente à \(N=15\), précritique à \(N=18\),
maximale à \(N=20\), puis devient plus globalement source à \(N=22\) et
\(N=24\).

Ces résultats suggèrent que les excitations BuP ne sont pas simplement des
perturbations localisées de \(W_{ij}\), mais des phases dynamiques internes du
réseau. Une excitation source peut être métastable, puis relaxer vers un
régime non-source. Cette observation ouvre une voie vers une interprétation de
la matière comme défaut métastable du réseau d'intrication, sans toutefois
identifier encore les particules du Modèle Standard.

\end{abstract}

\tableofcontents

% ─────────────────────────────────────────────────────────────
\section{Introduction}
% ─────────────────────────────────────────────────────────────

Le programme Bottom-Up Quantum Gravity (BuP) repose sur l'idée que l'espace, la
géométrie et la gravitation ne sont pas fondamentaux. Ils émergent de la structure
d'intrication d'un état quantique global. L'objet de base est une matrice
d'information mutuelle,
\[
W_{ij}=I(i:j),
\]
à partir de laquelle on reconstruit une distance d'intrication, une métrique
effective \(g_{\mu\nu}^{\rm ent}\), une dimension spectrale locale \(d_s(x)\) et
un couplage gravitationnel effectif.

Paper~5 a formulé une équation d'Einstein effective émergente,
\[
G_{\mu\nu}[g^{\rm ent}]
=
8\pi G_{\rm eff}(d_s)T_{\mu\nu}^{\rm matter}
+
8\pi G T_{\mu\nu}^{\rm ent}[d_s].
\]
Paper~6 a proposé une étape plus fondamentale :
\[
\mathcal{F}_{ij}[W]=0,
\]
ou encore
\[
\frac{\delta S_{\rm BuP}[W]}{\delta W_{ij}}=0.
\]
Dans cette formulation, la matière ne doit pas être ajoutée au réseau comme une
source externe. Elle doit émerger comme une excitation localisée de \(W_{ij}\).

Le présent papier explore la dynamique de ces excitations sous relaxation d'une
action candidate \(S_{\rm BuP}[W]\). La question centrale est :

\[
\boxed{
\text{Une excitation source reste-t-elle source sous relaxation de }S_{\rm BuP} ?
}
\]

La réponse numérique obtenue ici est non triviale. L'excitation ne conserve pas
nécessairement son caractère source. Elle peut changer de phase interne,
passer par un régime métastable, puis rejoindre un régime tardif non-source.

L'objectif de Paper~7 est donc de montrer que la matière émergente dans BuP ne
doit pas être comprise uniquement comme une bosse locale de \(\delta W\), mais
comme une phase dynamique de cette excitation dans le paysage variationnel de
\(S_{\rm BuP}\).

% ─────────────────────────────────────────────────────────────
\section{Rappel de Paper 6 : équation pure BuP}
% ─────────────────────────────────────────────────────────────

Paper~6 a introduit l'idée que l'équation fondamentale de BuP doit porter
directement sur le réseau d'information mutuelle :
\[
\mathcal{F}_{ij}[W]=0.
\]

Une écriture variationnelle possible est :
\[
\frac{\delta S_{\rm BuP}[W]}{\delta W_{ij}}=0.
\]

Dans cette perspective, l'équation d'Einstein effective de Paper~5 doit apparaître
comme une limite continue :
\[
\mathcal{F}_{ij}[W]=0
\quad
\Longrightarrow
\quad
G_{\mu\nu}[g^{\rm ent}]
=
8\pi G_{\rm eff}(d_s)T_{\mu\nu}^{\rm matter}
+
8\pi G T_{\mu\nu}^{\rm ent}[d_s].
\]

La chaîne d'émergence géométrique est :
\[
W_{ij}
\longrightarrow
d_{ij}^{\rm ent}
\longrightarrow
g_{\mu\nu}^{\rm ent}
\longrightarrow
d_s(x)
\longrightarrow
G_{\mu\nu}[g^{\rm ent}].
\]

La matière doit émerger d'une branche complémentaire :
\[
\delta W_{ij}^{\rm loc}
\longrightarrow
T_{\mu\nu}^{\rm matter}.
\]

Paper~7 se concentre précisément sur cette deuxième branche : la dynamique de
\(\delta W^{\rm loc}\).

% ─────────────────────────────────────────────────────────────
\section{Action candidate \(S_{\rm BuP}[W]\)}
% ─────────────────────────────────────────────────────────────

Nous considérons une action candidate de la forme générale :
\[
S_{\rm BuP}[W]
=
S_{\rm spec}[W]
+
S_{\rm geom}[W]
+
S_{\rm loc}[W]
+
S_{\rm norm}[W]
+
S_{\rm exc}[W].
\]

Le terme spectral contrôle la dimension spectrale :
\[
S_{\rm spec}[W]
=
\alpha
\left[
d_s(W)-d_\star
\right]^2.
\]

Le terme géométrique contrôle un proxy discret de courbure :
\[
S_{\rm geom}[W]
=
\beta
\left\langle
\left(G_{ij}^{\rm proxy}\right)^2
\right\rangle_{(i,j)}.
\]

Le terme de localité pénalise les connexions fortes à grande distance
d'intrication :
\[
S_{\rm loc}[W]
=
\gamma
\left\langle
W_{ij}\left(d_{ij}^{\rm ent}\right)^2
\right\rangle_{(i,j)}.
\]

Le terme de normalisation contrôle la masse totale d'intrication :
\[
S_{\rm norm}[W]
=
\mu
\left(
\frac{\sum_{i<j}W_{ij}-\sum_{i<j}W^{(0)}_{ij}}
{\sum_{i<j}W^{(0)}_{ij}}
\right)^2.
\]

Enfin, le terme d'excitation contrôle les déformations localisées :
\[
S_{\rm exc}[W]
=
m_{\rm exc}
\left[
\left\langle
(\delta W_{ij})^2
\right\rangle
+
\eta_{\rm grad}
S_{\rm smooth}[\delta W]
\right].
\]

Dans les simulations de ce papier, \(S_{\rm BuP}\) est utilisée comme action
candidate effective. Elle ne constitue pas encore l'action finale de BuP. Son
rôle est de fournir un protocole variationnel permettant d'étudier la relaxation
des excitations localisées.

% ─────────────────────────────────────────────────────────────
\section{Excitations localisées et observable matière-courbure}
% ─────────────────────────────────────────────────────────────

Une excitation localisée est construite sous la forme :
\[
W_{ij}
=
W_{ij}^{(0)}
+
\delta W_{ij}^{\rm loc}.
\]

Dans le protocole numérique, on injecte une perturbation gaussienne localisée :
\[
\delta W_{ij}^{\rm loc}
=
A f_i f_j,
\]
avec
\[
f_i
=
\exp
\left[
-\frac{d_{\rm ent}^2(i,i_0)}{2\sigma^2}
\right].
\]

Le proxy discret de matière est :
\[
T_{ij}^{\rm matter,proxy}
=
(\delta W_{ij})^2.
\]

La réponse géométrique est mesurée par :
\[
\delta G_{ij}^{\rm proxy}
=
G_{ij}^{\rm proxy}[W]
-
G_{ij}^{\rm proxy}[W^{(0)}],
\]
où
\[
G_{ij}^{\rm proxy}
=
-\kappa_{ij}
+
\frac12R_{\rm edge}.
\]

L'observable principale est :
\[
\mathcal{O}_{\rm matter}
=
\rho_{\rm Spearman}
\left(
T_{ij}^{\rm matter,proxy},
|\delta G_{ij}^{\rm proxy}|
\right).
\]

Si
\[
\mathcal{O}_{\rm matter}>0,
\]
l'excitation est interprétable comme source géométrique effective.
Si
\[
\mathcal{O}_{\rm matter}<0,
\]
la réponse devient anti-corrélée : l'excitation perd son caractère source.

% ─────────────────────────────────────────────────────────────
\section{Protocole numérique}
% ─────────────────────────────────────────────────────────────

Le protocole utilisé est le suivant :

\begin{enumerate}
  \item charger une matrice d'information mutuelle \(W^{(0)}(\lambda)\) ;
  \item construire le graphe \(k\)-NN associé ;
  \item calculer les distances d'intrication, la dimension spectrale et le proxy \(G_{ij}^{\rm proxy}\) ;
  \item injecter une excitation localisée \(\delta W^{\rm loc}\) ;
  \item mesurer \(\mathcal{O}_{\rm matter}^{\rm init}\) ;
  \item relaxer \(W\) selon une descente de \(S_{\rm BuP}\) ;
  \item mesurer \(\mathcal{O}_{\rm matter}^{\rm final}\) ;
  \item classifier le régime dynamique.
\end{enumerate}

La relaxation est réalisée sous la forme :
\[
W(t+1)
=
W(t)
-
\eta
\frac{\nabla_W S_{\rm BuP}}{\|\nabla_W S_{\rm BuP}\|}.
\]

Une procédure de line-search est utilisée pour n'accepter que les pas qui
diminuent l'action :
\[
S_{\rm BuP}[W(t+1)]
\leq
S_{\rm BuP}[W(t)].
\]

La dimension cible \(d_\star\) est prise en mode automatique :
\[
d_\star=d_s(W^{(0)}).
\]
Cela évite de forcer les petits graphes \(N=20\) vers une dimension cible
incompatible avec leur régime spectral effectif.

% ─────────────────────────────────────────────────────────────
\section{Résultat principal : transformation de phase sous relaxation}
% ─────────────────────────────────────────────────────────────

Nous définissons quatre classes dynamiques :

\[
(+,+):\quad
\mathcal{O}_{\rm init}>0,\quad
\mathcal{O}_{\rm final}>0,
\]
correspondant à une source stable.

\[
(+,-):\quad
\mathcal{O}_{\rm init}>0,\quad
\mathcal{O}_{\rm final}<0,
\]
correspondant à une source qui perd sa cohérence.

\[
(-,+):\quad
\mathcal{O}_{\rm init}<0,\quad
\mathcal{O}_{\rm final}>0,
\]
correspondant à une source qui émerge après relaxation.

\[
(-,-):\quad
\mathcal{O}_{\rm init}<0,\quad
\mathcal{O}_{\rm final}<0,
\]
correspondant à un régime non-source.

Dans le scan de référence à \(20\) pas de relaxation, on obtient :
\[
\lambda_{\rm init}\simeq0.265,
\]
pour le seuil initial, et
\[
\lambda_{\rm final}\simeq0.876,
\]
pour le seuil final.

Les résultats montrent que plusieurs excitations initialement source
deviennent anti-corrélées après relaxation. C'est le cas des régimes
\[
\lambda=0.29,\quad0.43,\quad0.57.
\]
Ces régimes sont donc classés comme \emph{source décohérée}.

Cela montre que le caractère source d'une excitation n'est pas seulement
une propriété statique de l'injection initiale. Il dépend de la dynamique
variationnelle de \(S_{\rm BuP}\).

\begin{figure}[htbp]
\centering
\includegraphics[width=0.92\textwidth]{../figures/fig_phase_scan_lambda.png}
\caption{
Évolution de \(\mathcal{O}_{\rm matter}\) avant et après relaxation de
\(S_{\rm BuP}\). Le seuil initial se situe autour de
\(\lambda_{\rm init}\simeq0.265\), tandis que le seuil final à relaxation
intermédiaire se situe autour de \(\lambda_{\rm final}\simeq0.876\).
}
\label{fig:phase-scan}
\end{figure}

% ─────────────────────────────────────────────────────────────
\section{Dépendance au temps de relaxation}
% ─────────────────────────────────────────────────────────────

Nous avons étudié la dépendance au nombre de pas de relaxation pour le cas
\(N=20\). Le seuil initial reste constant :
\[
\lambda_{\rm init}\simeq0.265.
\]

En revanche, le seuil final dépend fortement du temps de relaxation.

\begin{table}[htbp]
\centering
\begin{tabular}{c c c}
\toprule
Nombre de pas & Seuil final & Interprétation \\
\midrule
5 & \(0.215\) & régime transitoire court \\
10--20 & \(0.876\) & plateau métastable source \\
40--320 & \(0.280\) & régime tardif non-source \\
\bottomrule
\end{tabular}
\caption{
Évolution du seuil source/non-source en fonction du nombre de pas de
relaxation pour \(N=20\). Le plateau intermédiaire autour de \(0.876\) est
interprété comme un régime métastable source. Le régime autour de \(0.280\)
persiste jusqu'à \(320\) pas et constitue un attracteur tardif candidat.
}
\label{tab:steps-crossing}
\end{table}

La structure observée est donc :
\[
\lambda_{\rm init}\simeq0.265,
\]
\[
\lambda_{\rm meta}\simeq0.876,
\]
\[
\lambda_{\rm late}\simeq0.280.
\]

Le régime \(\lambda_{\rm meta}\simeq0.876\) apparaît pour des temps
intermédiaires, entre \(10\) et \(20\) pas. Le régime
\(\lambda_{\rm late}\simeq0.280\) apparaît à partir de \(40\) pas et reste stable
jusqu'à \(320\) pas.

Cette séparation suggère que la relaxation de \(S_{\rm BuP}[W]\) ne produit pas
une évolution monotone simple. Elle révèle au moins deux régimes dynamiques :
un plateau métastable source et un régime tardif non-source.

\begin{figure}[htbp]
\centering
\includegraphics[width=0.85\textwidth]{../figures/fig_crossing_vs_steps.png}
\caption{
Évolution du seuil final en fonction du nombre de pas de relaxation.
Le régime intermédiaire autour de \(0.876\) correspond à un plateau métastable,
tandis que le régime tardif autour de \(0.280\) reste stable jusqu'à \(320\)
pas.
}
\label{fig:crossing-steps}
\end{figure}

% ─────────────────────────────────────────────────────────────
\section{Tendance en taille finie et transition de phase en \(N\)}
% ─────────────────────────────────────────────────────────────

Nous avons ensuite testé la dépendance en taille du réseau pour
\[
N=15,\;18,\;20,\;22,\;24.
\]
Le but est de vérifier si la dynamique métastable observée à \(N=20\) est déjà
présente dans des réseaux plus petits, et comment elle évolue lorsque \(N\)
augmente.

\begin{table}[htbp]
\centering
\begin{tabular}{c c c c l}
\toprule
\(N\) & Steps & \(\lambda_{\rm init}\) & \(\lambda_{\rm final}\) & Interprétation \\
\midrule
15 & 20  & \(0.042\) & \(0.042\) & pas de dynamique visible \\
15 & 160 & \(0.042\) & \(0.042\) & pas de dynamique visible \\
18 & 20  & \(0.522\) & \(0.478\) & régime précritique \\
18 & 160 & \(0.522\) & \(0.478\) & régime précritique stable \\
20 & 20  & \(0.265\) & \(0.876\) & plateau métastable fort \\
20 & 160 & \(0.265\) & \(0.280\) & régime tardif non-source \\
22 & 20  & — & \(0.359\) & source initial dominant \\
22 & 160 & — & \(0.360\) & seuil dynamique stable \\
24 & 20  & \(0.427\) & — & phase source globale \\
24 & 160 & \(0.427\) & \(0.521\) & sélection tardive \\
\bottomrule
\end{tabular}
\caption{
Tendance en taille finie des seuils source/non-source.
Le symbole « — » indique l’absence de seuil initial (\(\mathcal{O}_{\rm init}>0\)
sur toute la grille testée) ou l’absence de seuil final (\(\mathcal{O}_{\rm final}>0\)
sur toute la grille). La dynamique est absente à \(N=15\), précritique à \(N=18\),
fortement métastable à \(N=20\), puis devient plus globalement source à \(N=22\)
et \(N=24\). Le cas \(N=24\) montre une phase source complète à \(20\) pas,
suivie d'une sélection tardive à \(160\) pas.
}
\label{tab:finite-size}
\end{table}

La comparaison montre une structure non monotone mais cohérente.

À \(N=15\), la relaxation ne modifie pas le seuil :
\[
\lambda_{\rm init}
=
\lambda_{\rm final}
\simeq0.042.
\]
La dynamique interne des excitations n'est donc pas encore visible.

À \(N=18\), une dynamique précritique apparaît :
\[
\lambda_{\rm init}\simeq0.522,
\qquad
\lambda_{\rm final}\simeq0.478.
\]
Le seuil est modifié par la relaxation, mais il reste stable entre \(20\) et
\(160\) pas. Le plateau métastable observé à \(N=20\) n'est pas encore présent.

À \(N=20\), le système présente la structure la plus riche :
\[
\lambda_{\rm final}(20)\simeq0.876,
\]
puis
\[
\lambda_{\rm final}(160)\simeq0.280.
\]
Nous interprétons ce cas comme un régime critique ou quasi-critique où la
métastabilité source devient clairement visible. La forte séparation entre
le plateau métastable et le régime tardif, ainsi que la sensibilité au temps
de relaxation, plaident en faveur de ce caractère critique.

À \(N=22\), le seuil initial disparaît car
\[
\mathcal{O}_{\rm matter}^{\rm init}>0
\]
sur toute la grille de \(\lambda\) testée. Les excitations injectées sont donc
source sur toute la plage explorée. Après relaxation, un seuil dynamique
stable apparaît autour de
\[
\lambda_{\rm final}^{N=22}\simeq0.36.
\]

À \(N=24\), la relaxation intermédiaire rend toutes les excitations source :
\[
\lambda_{\rm final}(20)=\text{—}\ (\mathcal{O}_{\rm final}>0\ \text{partout}).
\]
À relaxation longue, un seuil réapparaît :
\[
\lambda_{\rm final}(160)\simeq0.521.
\]

Ces résultats suggèrent que la matière effective devient une propriété
collective du réseau lorsque \(N\) augmente. Le cas \(N=20\) semble marquer une
région critique de taille finie, tandis que \(N=22\) et \(N=24\) indiquent une
phase plus globalement source.

Nous n'interprétons pas encore cette tendance comme une limite asymptotique
\(N\to\infty\). Les données montrent une dépendance non monotone en \(N\), qui
devra être explorée par des tailles supplémentaires.

% ─────────────────────────────────────────────────────────────
\section{Robustesse vis-à-vis des coefficients de l'action}
% ─────────────────────────────────────────────────────────────

Nous avons testé la dépendance du seuil final aux coefficients de l'action
candidate \(S_{\rm BuP}\).

Le seuil métastable obtenu à \(10\) à \(20\) pas reste stable lorsque l'on varie
séparément :
\[
\gamma=0.005,\;0.01,\;0.02,\;0.05,
\]
\[
\beta=0.02,\;0.05,\;0.1,\;0.2,\;0.5,
\]
et
\[
m_{\rm exc}=0,\;0.01,\;0.05,\;0.1,\;0.2.
\]

Dans tous ces scans, le seuil final reste proche de :
\[
\lambda_{\rm final}\simeq0.876.
\]

Cette robustesse suggère que le plateau métastable n'est pas un artefact d'un
réglage fin des coefficients de l'action. Il semble plutôt refléter une propriété
de la famille de réseaux \(W_{ij}(\lambda)\) sous relaxation variationnelle.

Cette conclusion doit toutefois être nuancée : le protocole utilise un gradient
normalisé et une line-search, ce qui peut réduire la sensibilité aux coefficients.
La robustesse observée doit donc être comprise comme une robustesse dans le
protocole numérique considéré, et non comme une preuve d'universalité absolue.

% ─────────────────────────────────────────────────────────────
\section{Structure interne des excitations}
% ─────────────────────────────────────────────────────────────

Les simulations montrent que la masse informationnelle et le rayon d'une
excitation ne suffisent pas à déterminer son comportement géométrique.

Dans plusieurs runs, on observe :
\[
M_{\rm exc}^{\rm init}
\simeq
M_{\rm exc}^{\rm final},
\]
\[
R_{\rm exc}^{\rm init}
\simeq
R_{\rm exc}^{\rm final},
\]
tandis que :
\[
\mathcal{O}_{\rm matter}^{\rm init}
\neq
\mathcal{O}_{\rm matter}^{\rm final}.
\]

Cela signifie que l'excitation peut conserver sa masse et sa taille tout en
changeant sa cohérence géométrique interne. La nature d'une excitation BuP ne
dépend donc pas seulement de quantités globales comme \(M_{\rm exc}\) ou
\(R_{\rm exc}\), mais aussi de sa structure interne.

Les candidats naturels pour caractériser cette structure sont :
\[
{\rm Spec}(L_{\rm exc}),
\]
le spectre du Laplacien local de l'excitation,
\[
M_4
=
\sum_i r_i^4\rho_i^{\rm exc},
\]
un moment de forme,
ou encore des invariants topologiques tels que les cycles excités du sous-graphe
porteur de \(\delta W^{\rm loc}\).

Dans une version plus complète de BuP, une particule pourrait correspondre à une
classe stable de ces invariants internes.

% ─────────────────────────────────────────────────────────────
\section{Discussion : matière métastable}
% ─────────────────────────────────────────────────────────────

Les résultats de Paper~7 suggèrent une lecture programmatique :

\[
\boxed{
\text{la matière BuP pourrait être une phase métastable de l'intrication.}
}
\]

Dans cette perspective, une particule n'est pas un objet fondamental ajouté au
réseau. Elle est un défaut organisé de \(W_{ij}\), capable de produire une
réponse géométrique source pendant un temps effectif fini.

Une durée de vie effective peut être définie par :
\[
\tau_{\rm life}
=
\min
\left\{
t:
\mathcal{O}_{\rm matter}(t)\leq0
\right\}.
\]

Une excitation est interprétable comme particule effective tant que :
\[
\mathcal{O}_{\rm matter}(t)>0.
\]

La stabilité d'une particule BuP pourrait alors résulter de trois mécanismes :

\begin{enumerate}
  \item une barrière d'action élevée dans \(S_{\rm BuP}\) ;
  \item une protection topologique ou spectrale ;
  \item une symétrie interne empêchant la relaxation vers le régime non-source.
\end{enumerate}

Cette interprétation reste spéculative. Les pas de relaxation utilisés ici ne
sont pas encore identifiés au temps physique. Ils doivent être compris comme un
temps effectif de relaxation du réseau. Le lien avec le temps cosmologique ou le
temps émergent BuP reste à établir.

Néanmoins, les résultats montrent que le caractère source de la matière
émergente n'est pas trivialement stable. Il dépend de la trajectoire de
relaxation dans le paysage de \(S_{\rm BuP}\).

% ─────────────────────────────────────────────────────────────
\section{Limites et perspectives}
% ─────────────────────────────────────────────────────────────

Cette étude reste préliminaire.

Premièrement, l'action \(S_{\rm BuP}\) utilisée ici est une action candidate.
Elle n'est pas encore dérivée de principes premiers.

Deuxièmement, les matrices \(W_{ij}\) utilisées ne sont pas générées comme
solutions exactes de
\[
\frac{\delta S_{\rm BuP}}{\delta W_{ij}}=0.
\]
Elles proviennent d'une famille de matrices MI déjà construites, puis relaxées.

Troisièmement, le temps de relaxation n'est pas encore identifié au temps
physique.

Quatrièmement, les excitations injectées sont des perturbations gaussiennes
simples. Des excitations topologiques, orientées ou complexes devront être
testées.

Les perspectives immédiates sont :

\begin{enumerate}
  \item tester plusieurs centres d'excitation ;
  \item tester plusieurs amplitudes et largeurs \(\sigma\) ;
  \item introduire des invariants topologiques de boucle ;
  \item ajouter un terme de protection \(S_{\rm protect}\) dans l'action ;
  \item tester si une excitation peut rester source jusqu'à \(320\) pas ou plus ;
  \item construire un diagramme de phase \((\lambda,t)\) complet.
\end{enumerate}

% ─────────────────────────────────────────────────────────────
\section{Conclusion}
% ─────────────────────────────────────────────────────────────

Paper~7 explore la dynamique variationnelle des excitations BuP sous relaxation
d'une action candidate \(S_{\rm BuP}[W]\). Cette dynamique est une manière
d’explorer le paysage de l’équation pure \(\mathcal{F}_{ij}[W]=0\) (Paper~6).

Le résultat principal est que le caractère source d'une excitation
localisée n'est pas conservé automatiquement. L'observable
\[
\mathcal{O}_{\rm matter}
=
\rho_{\rm Spearman}
\left(
T_{ij}^{\rm matter,proxy},
|\delta G_{ij}^{\rm proxy}|
\right)
\]
évolue sous relaxation et révèle plusieurs régimes dynamiques.

Pour \(N=20\), nous identifions :
\[
\lambda_{\rm init}\simeq0.265,
\]
comme seuil initial,
\[
\lambda_{\rm meta}\simeq0.876,
\]
comme plateau métastable source,
et
\[
\lambda_{\rm late}\simeq0.280,
\]
comme régime tardif non-source, stable jusqu'à \(320\) pas.

Le scan en taille finie montre que cette dynamique dépend fortement de \(N\).
À \(N=15\), elle est absente. À \(N=18\), un régime précritique apparaît. À
\(N=20\), la métastabilité source est maximale. À \(N=22\) et \(N=24\),
les excitations deviennent plus globalement source, avec des seuils
dynamiques tardifs.

Ces résultats suggèrent que la matière émergente dans BuP ne doit pas être
comprise seulement comme une excitation localisée de \(W_{ij}\), mais comme une
phase dynamique interne de cette excitation dans le paysage variationnel de
\(S_{\rm BuP}\).

La conclusion conceptuelle prudente est :
\[
\boxed{
\text{une particule BuP pourrait correspondre à un défaut métastable source du réseau d'intrication.}
}
\]

Sa stabilité demanderait alors un mécanisme de protection : barrière d'action,
invariant topologique, invariant spectral ou symétrie interne. Sans protection,
une excitation source peut relaxer vers un régime tardif non-source.

Le présent papier ne prétend pas identifier les particules du Modèle Standard.
Il établit plutôt un résultat intermédiaire : les excitations BuP possèdent une
dynamique interne non triviale sous relaxation de \(S_{\rm BuP}\), avec
métastabilité, régimes tardifs et effets de taille finie.

% ─────────────────────────────────────────────────────────────
\begin{thebibliography}{9}
% ─────────────────────────────────────────────────────────────

\bibitem{hamdad-paper5}
F.~Hamdad,
\textit{Équation d'Einstein effective émergente dans la Gravité Quantique Bottom-Up},
BuP Paper 5, GitHub repository, 2026.

\bibitem{hamdad-paper6}
F.~Hamdad,
\textit{Équation pure BuP : de l'intrication à l'équation d'Einstein effective},
BuP Paper 6, GitHub repository, 2026.

\bibitem{ollivier}
Y.~Ollivier,
\textit{Ricci curvature of Markov chains on metric spaces},
Journal of Functional Analysis, 256, 810--864, 2009.

\bibitem{spectral-dimension}
D.~Benedetti,
\textit{Fractal properties of quantum spacetime},
Physical Review Letters, 102, 111303, 2009.

\end{thebibliography}

\end{document}
